% titel: Blatt 4, A5: Photoelektrischer Effekt an Molybdän
% skript: Mortimer Kap. 6 (Photoeffekt, Austrittsarbeit)

% Teil I: Werkzeuge, kein Lösungsweg.

Beim Photoeffekt gibt ein Photon seine Energie an ein Metall-Elektron ab;
nur oberhalb der (materialabhängigen) Frequenz $\nu_0$ kommt ein Elektron
heraus --- die Energie über der Austrittsarbeit erscheint als kinetische
Energie.

\begin{werkzeugkasten}[Photoeffektgleichung]
\begin{equation}
  E_{\text{photon}} = E_{\text{kin}} + W,
  \qquad
  E_{\text{photon}} = h\,\nu = \frac{h\,c}{\lambda},
  \qquad
  W = h\,\nu_0.
\end{equation}
$\nu_0$ = Grenzfrequenz (kleinste Frequenz, die Emission ermöglicht),
$W$ = Austrittsarbeit, $E_{\text{kin}}$ = maximale kinetische Energie des
austretenden Elektrons.

\textit{Skript:} Mortimer Kap. 6 (Photoelektrischer Effekt), Übung 4/5.
\end{werkzeugkasten}

\begin{achtung}[Darauf achten]
\begin{itemize}
  \item Mindestenergie = Austrittsarbeit $W = h\,\nu_0$ (bei gerader Frequenz
        $= 1.09\times10^{15}$, direkt aus der Grenzfrequenz).
  \item Wellenlänge aus Frequenz: $\lambda = c/\nu_0$; nicht
        $E_{\text{photon}}$ mit $W$ verwechseln.
  \item $E_{\text{kin}} = E_{\text{photon}} - W$; beide Energien vorher in
        J (oder beide in eV) umrechnen.
  \item Einheit $s^{-1}$ bzw.\ Hz für $\nu_0$, Wellenlänge in m (120 nm
        $= \qty{1.20e-7}{m}$).
\end{itemize}
\end{achtung}
